% =========================================================
% 12. default_get() ORM Method
% =========================================================
\section{\texttt{default\_get()} ORM Method}

\begin{frame}[standout]
  \texttt{default\_get()} ORM Method
\end{frame}

\begin{frame}[fragile]{\texttt{DEFAULT\_GET()} METHOD}
  \begin{itemize}
    \item The \texttt{default\_get()} method returns the \textbf{default values}
          used when creating a new record in the UI or through the ORM.
    \item It collects defaults from field definitions, context, and company settings.
    \item The method returns a dictionary of field values.
  \end{itemize}
  \begin{block}{Method Syntax}
\begin{lstlisting}
@api.model
def default_get(self, fields_list):
    return super().default_get(fields_list)
\end{lstlisting}
    Important parts: \texttt{fields\_list}.
  \end{block}
\end{frame}

\begin{frame}[fragile]{BREAKING DOWN THE METHOD SIGNATURE}
  \begin{itemize}
    \item \texttt{@api.model}: Model-level method.
    \item \texttt{self}: The model requesting defaults.
    \item \texttt{fields\_list}: List of field names for which defaults are needed.
  \end{itemize}
\begin{lstlisting}
fields_list = ['name', 'gender', 'active', 'company_id']
\end{lstlisting}
  \begin{itemize}
    \item The UI calls this when you click \textbf{New} on a form.
  \end{itemize}
\end{frame}

\begin{frame}[fragile]{WHAT DOES \texttt{DEFAULT\_GET()} RETURN?}
  \begin{itemize}
    \item It returns a dictionary of default values.
  \end{itemize}
\begin{lstlisting}
defaults = self.env['student.student'].default_get([
    'name', 'gender', 'active',
])
print(defaults)
\end{lstlisting}
  Output is (example):
\begin{lstlisting}
{'active': True, 'gender': 'male'}
\end{lstlisting}
  \begin{itemize}
    \item Only fields that have defaults appear in the dictionary.
    \item Missing keys mean ``no default''.
  \end{itemize}
\end{frame}

\begin{frame}[fragile]{WHERE DEFAULTS COME FROM}
  \begin{enumerate}
    \item Field-level default:
\begin{lstlisting}
active = fields.Boolean(default=True)
gender = fields.Selection(..., default='male')
\end{lstlisting}
    \item Context default:
\begin{lstlisting}
self.env['student.student'].with_context(
    default_gender='female'
).default_get(['gender'])
\end{lstlisting}
    \item Custom logic inside an overridden \texttt{default\_get()}.
  \end{enumerate}
\end{frame}

\begin{frame}[fragile]{OVERRIDE PATTERN}
\begin{lstlisting}
@api.model
def default_get(self, fields_list):
    defaults = super().default_get(fields_list)
    if 'name' in fields_list:
        defaults['name'] = 'New Student'
    if 'active' in fields_list:
        defaults['active'] = True
    return defaults
\end{lstlisting}
  \begin{itemize}
    \item Always call \texttt{super()} first, then adjust the dictionary.
    \item Only set keys that were requested in \texttt{fields\_list}.
  \end{itemize}
\end{frame}

\begin{frame}[fragile]{METHOD CALL}
  \textbf{Method Call}
\begin{lstlisting}
vals = self.env['student.student'].default_get([
    'name', 'birth_date', 'gender', 'active',
])
\end{lstlisting}
  \begin{itemize}
    \item Here, \texttt{.default\_get(...)} is the method call.
    \item Return type reminder: \textbf{dict}.
  \end{itemize}
\end{frame}
